\paragraph{A practical $3+2$ reduction on the pseudoinverse slice.}

The full spectral representation of the reduced quartic tensor is
\begin{equation}
\Tau' = R(\alpha,\beta,\gamma)\,
\mathrm{diag}(\lambda_1,\lambda_2,\lambda_3)\,
R(\alpha,\beta,\gamma)^T,
\label{eq:tauprime_spectral_3plus3}
\end{equation}
which provides a natural $3+3$ description in terms of three eigenvalues and
three orientation angles.  This is the fundamental tensorial
parametrization of the quartic problem.

However, when the tensor is reconstructed from Watson constants by means of the
Moore--Penrose pseudoinverse, one does not work with an arbitrary element of
the spectral manifold, but with the specific representative selected by the
pseudoinverse gauge fixing.  In the notation introduced above, this
representative satisfies
\begin{equation}
F(\Tau')=
2\Tau'_{ab}+(\Sigma-1)\Tau'_{ac}-(\Sigma+1)\Tau'_{bc}=0,
\label{eq:tauprime_pinv_slice}
\end{equation}
with
\begin{equation}
\Sigma=\frac{2A-B-C}{B-C}.
\end{equation}

Equation~\eqref{eq:tauprime_pinv_slice} removes one coordinate from the full
$3+3$ spectral representation and therefore defines a practical $3+2$
description of the pseudoinverse representative.  A convenient choice is to
take
\begin{equation}
(\lambda_1,\lambda_2,\lambda_3,\alpha,\beta)
\end{equation}
as independent coordinates and to determine the remaining Euler angle $\gamma$
from the scalar gauge condition
\begin{equation}
F\!\left(
R(\alpha,\beta,\gamma)\,
\mathrm{diag}(\lambda_1,\lambda_2,\lambda_3)\,
R(\alpha,\beta,\gamma)^T
\right)=0.
\label{eq:gamma_from_gauge}
\end{equation}

This $3+2$ representation should not be interpreted as a global
representation-independent parametrization of all quartic tensors.  Rather, it
is a coordinate chart on the five-dimensional pseudoinverse slice selected in
the inverse reconstruction problem.  In this sense, the $3+3$ spectral form is
the fundamental tensorial description, whereas the $3+2$ reduction is a
practical parametrization of the specific quartic tensor representative adopted
in the present implementation.
